
Uma característica marcante de linguagens de programação funcionais é que seus programas costumam se assemelhar muito mais a definições matemáticas do que em linguagens imperativas, onde os programas se assemelham a uma sequência de passos que a computador deve seguir para calcular o desejado. Em vez de descrever uma sequência de instruções que modifica o estado da memória, um programa funcional especifica diretamente a relação entre seus argumentos e seu resultado.

Para a definição recursiva em Haskell do produto interno, por exemplo, sua implementação é praticamente uma transcrição literal da definição matemática. Cada equação da definição matemática corresponde a uma equação do programa. Em particular, os padrões utilizados na definição da função desempenham exatamente o papel da distinção de casos feita na definição recursiva.

Essa semelhança entre definições matemáticas e programas não ocorre por acaso. Em Haskell, funções são livres de efeitos colaterais: seu resultado depende exclusivamente dos argumentos recebidos, sem modificar qualquer estado externo. Consequentemente, uma chamada de função pode ser encarada simplesmente como uma expressão matemática cujo valor é determinado pelas equações que definem essa função.

Essa propriedade permite tratar programas da mesma maneira que tratamos expressões algébricas. Assim como é possível substituir uma expressão matemática por outra que lhe seja igual, também é possível substituir uma chamada de função pela expressão correspondente dada por sua definição. Esse princípio recebe o nome de raciocínio equacional e constitui uma das principais ferramentas para demonstrar propriedades e equivalência entre programas funcionais.

Com essa perspectiva, demonstrações de igualdade entre programas podem ser conduzidas de maneira análoga às demonstrações de igualdade entre expressões matemáticas. Em muitos casos, provas por indução são suficientes para estabelecer as propriedades desejadas. Entretanto, existem outras estratégias de demonstração que exploram diretamente as equações que definem os programas, como discutido em \cite{Gibbons2005}. Neste trabalho, adotaremos essa segunda abordagem. Em particular, reescreveremos os algoritmos clássico e modificado de Gram-Schmidt em uma linguagem baseada em streams e demonstraremos sua equivalência por meio de raciocínio equacional, sem recorrer a uma prova por indução.
